\chapter{Problem Statement and Objectives}

\section{Problem Statement}
The aviation industry faces a persistent challenge in ensuring that passengers are adequately informed about safety procedures. Traditional safety briefings are passive and often ignored. In the event of an emergency, this lack of engagement can lead to panic and a failure to follow critical instructions. The problem is to design a training system that overcomes this passivity and ensures that passengers actively learn and retain safety information.

\section{Objectives}
The primary objectives of this project are:
\begin{itemize}
    \item \textbf{To Develop a Realistic VR Environment:} Create a 3D simulation of an aircraft cabin (Boeing 737) to provide context and immersion.
    \item \textbf{To Implement Interactive Learning Modules:} Replace passive watching with active participation using a custom-built video player and spatial UI.
    \item \textbf{To Verify User Understanding:} Design and integrate a comprehensive assessment system (Quiz) that provides immediate feedback on safety knowledge.
    \item \textbf{To Optimize for Performance:} Utilize the lightweight StereoKit framework to ensure high frame rates and accessibility on standalone VR hardware (Meta Quest 3).
\end{itemize}

\section{Scope}
The scope of this project includes the development of a VR application that simulates a pre-flight safety briefing. It covers the following modules:
\begin{enumerate}
    \item Welcome and orientation.
    \item Video-based instruction on safety features (seatbelts, oxygen masks, life vests).
    \item Interactive quiz to test knowledge.
    \item Basic movement and navigation within the virtual cabin.
\end{enumerate}
The project is designed for the Meta Quest 3 headset but also supports a desktop fallback mode for testing.
